\par\noindent\textbf{(1)} 在 \(\alpha\) 固定时对 \(\lambda\) 求导，得 \(d\cos\beta\,d\beta/d\lambda=m\)，所以 \(d\beta/d\lambda=m/[d\cos\beta]\)。级次 \(m\) 的符号决定色散方向，\(\cos\beta\to0\) 时角度色散剧增。
\par\noindent\textbf{(2)} 出射波在固定基准坐标中的横向波矢 \(k_x(\omega)=(\omega/c)\sin\beta(\omega)\)。因此 \(dk_x/d\omega=c^{-1}[\sin\beta+\omega\cos\beta\,d\beta/d\omega]\)。由 \(\lambda=2\pi c/\omega\) 得 \(d\beta/d\omega=-m\lambda/(d\omega\cos\beta)\)，从而 \(dk_x/d\omega=c^{-1}(\sin\beta-m\lambda/d)=-\sin\alpha/c\)。
\par\noindent\textbf{(3)} 对复光场 \(E\propto\int A(\omega)e^{i[k_x(\omega)x+k_z(\omega)z-\omega t]}d\omega\)，群延迟等相位条件 \(t=x\,dk_x/d\omega+z\,dk_z/d\omega\)。故同一脉冲面的时间与横向位置耦合，局部脉冲前倾由群延迟梯度 \(\partial t/\partial x=dk_x/d\omega\)（在固定 \(z\)）定义，其数值依所选坐标和光栅安装几何。
\par\noindent\textbf{(4)} 靠近出射掠射 \(|\cos\beta|\ll1\)，小带宽可能引起大角散、光束截断与偏振效率变化；若级次不满足 \(|m\lambda/d-\sin\alpha|\le1\)，它不再是远场传播级次，不能对不存在的出射角继续求导。
